\section{Limitations}

When asking participants to name colors, aiming to detect differences in these labeling patterns, it is key to actually present the exactly same colors. In e.g. the WCS and \textcite{lindseyColor2014} this was ensured by the use of the 330 physical Musell chips of high quality.
For this experiment this was not possible due to the fact that we use an online experiment approach. Therefore, we had to accommodate to the fact that every participant uses their own potentially calibrated screen to present the colors and then name it. In fact, different screen models, even with different calibration settings, display colors significantly differently \parencite{kwak2000characterisation}. As a consequence, this noise is included into our data, compared to the in-lab settings.
However, there are techniques to counteract this issue such as gamma correction. \textcite{xiao2011visual} invented a procedure to determine the current state of calibration of a screen based on the observers’ visual luminance judgements for the three channels red, green and blue. This is a very promising approach that was also implemented as a pre screen test for PsyNet. However, testing this procedure and establishing the scientific justification exceeded the time available for this master's thesis.

As mentioned in the introduction, online experiments also impose other limitations. The most basic criticism is the inability to maintain controlled environmental conditions, a standard rigorously upheld in laboratory-based experiments. It cannot be prevented that participants multitask at the same time. While this is certainly true, the collected data showed high similarity to the in-lab experiment conducted by \textcite{lindseyColor2014}. An option to further justify the results of this experiment would be to do a control experiment under lab conditions. Comparison between the remote procedure and the control experiment could then reveal the extend of influence of the online conditions have.

The sample of 25 languages is already a substantial set but does not include enough languages of the global South in particular, South America and Africa \parencite{barrett2020towards, blasi2022over}. Nonetheless, the innovative recruitment strategy has yielded reliable data, breaking ground for future data collection efforts in these regions. Looking forward we can potentially recruit even larger group of languages and cultures.

Ultimately, we believe this work paves the way for using diverse recruiting and machine learning to understand the human mind.
The results show high consistency across Lucid and the more traditional recruiter Prolific.
This lays the foundation for further leveraging Lucid's enormous recruiting potential.
For many languages we saw newly emerging color terms which can in part be condensed into major color categories such as light blue.
This indicates that the set of BCT is not ultimately restricted to eleven terms.
It is the duty of future research to monitor this development.
Furthermore, the preserved cross language differences within the GPT4 data has proven LLMs to be able to contribute to the science of color terms and thus to decipher the relation between thought and language.
